% Part: first-order-logic
% Chapter: syntax-and-semantics
% Section: substitution

\documentclass[../../../include/open-logic-section]{subfiles}

\begin{document}

\olfileid{fol}{syn}{ext}

\olsection{ব্যাপ্তিগততা}

\begin{explain}
ব্যাপ্তিগততা, যাকে কখনও প্রাসঙ্গিকতাও বলা হয়, অনানুষ্ঠানিকভাবে এভাবে
প্রকাশ করা যায়: !!a{structure}~$\Struct M$-এ !!a{variable} আরোপ~$s$-এর
সাপেক্ষে !!{formula}~$!A$-এর পরিতৃপ্তির উপর প্রভাব ফেলে শুধু
!!{domain}-এর আকার এবং $!A$-তে বাস্তবে দেখা দেওয়া ভাষার উপাদানগুলিকে
$\Struct M$ ও~$s$ যে আরোপ দেয়।

ব্যাপ্তিগততার একটি তাৎক্ষণিক ফল হলো: দুটি !!{structure}s~$\Struct M$
ও~$\Struct M'$ যদি বাক্য~$!A$-তে দেখা দেওয়া ভাষার সব উপাদানে একমত হয়
এবং তাদের সংজ্ঞাক্ষেত্র একই হয়, তবে $!A$ সত্য কি না সে বিষয়েও
$\Struct M$ ও~$\Struct M'$-কে একমত হতে হবে।
\end{explain}

\begin{prop}[ব্যাপ্তিগততা]
  \ollabel{prop:extensionality}
  ধরি $!A$ একটি !!a{formula}; $\Struct M_1$ ও $\Struct M_2$ এমন দুটি
  !!{structure}s যাতে $\Domain{M_1} = \Domain{M_2}$; আর $s$ হলো
  $\Domain{M_1} = \Domain{M_2}$-এর উপর একটি চলরাশি-আরোপ। $!A$-তে
  সংঘটিত প্রত্যেক !!{constant}~$c$, সম্বন্ধসংকেত~$R$ এবং
  !!{function}~$f$-এর জন্য যদি $\Assign{c}{M_1} = \Assign{c}{M_2}$,
  $\Assign{R}{M_1}=\Assign{R}{M_2}$ এবং
  $\Assign{f}{M_1} = \Assign{f}{M_2}$ হয়, তাহলে
  $\Sat{M_1}{!A}[s]$ তখনই, যখন $\Sat{M_2}{!A}[s]$।
\end{prop}

\begin{proof}
  প্রথমে $t$-এর উপর আরোহ করে প্রমাণ করুন যে প্রত্যেক পদের জন্য
  $\Value{t}{M_1}[s] = \Value{t}{M_2}[s]$। তারপর এইমাত্র প্রমাণিত
  দাবিটি আরোহের ভিত্তিতে (যেখানে $!A$ পরমাণু) ব্যবহার করে $!A$-এর উপর
  আরোহের সাহায্যে প্রতিজ্ঞাটি প্রমাণ করুন।
\end{proof}

\begin{prob}
\olref[fol][syn][ext]{prop:extensionality}-এর প্রমাণটি বিস্তারিতভাবে
সম্পন্ন করুন।
\end{prob}

\begin{cor}[!!^{sentence}s-এর জন্য ব্যাপ্তিগততা]
  \ollabel{cor:extensionality-sent}
  ধরি $!A$ একটি !!a{sentence} এবং $\Struct{M_1}$, $\Struct{M_2}$
  \olref{prop:extensionality}-এর মতো। তাহলে $\Sat{M_1}{!A}$ তখনই,
  যখন $\Sat{M_2}{!A}$।
\end{cor}

\begin{proof}
\olref[ass]{cor:sat-sentence} ব্যবহার করে
\olref{prop:extensionality} থেকে ফলটি পাওয়া যায়।
\end{proof}

তদুপরি, পদের মান এবং কোনো !!a{structure} কোনো !!a{formula}-কে পরিতৃপ্ত
করে কি না, তা শুধু তার উপপদগুলির মানের উপর নির্ভর করে।

\begin{prop}\ollabel{prop:ext-terms}
ধরি $\Struct M$ একটি !!a{structure}, $t$ ও $t'$ দুটি পদ এবং $s$ একটি
চলরাশি-আরোপ। তাহলে $\Value{\Subst{t}{t'}{x}}{M}[s] =
\Value{t}{M}[\Subst{s}{\Value{t'}{M}[s]}{x}]$।
\end{prop}

\begin{proof}
$t$-এর উপর আরোহ করে।
\begin{enumerate}
\item $t$ ধ্রুবক হলে, ধরি $t\ident c$। তখন $\Subst{t}{t'}{x} = c$,
  এবং $\Value{c}{M}[s] = \Assign{c}{M} =
  \Value{c}{M}[\Subst{s}{\Value{t'}{M}[s]}{x}]$।

\item $t$,~$x$ ছাড়া অন্য চলরাশি হলে, ধরি $t \ident y$। তখন
  $\Subst{t}{t'}{x} = y$, এবং
  $\varAssign{s}{\Subst{s}{\Value{t'}{M}[s]}{x}}{x}$ বলে
  $\Value{y}{M}[s] =
  \Value{y}{M}[\Subst{s}{\Value{t'}{M}[s]}{x}]$।

\item $t \ident x$ হলে $\Subst{t}{t'}{x} = t'$। কিন্তু
  $\Subst{s}{\Value{t'}{M}[s]}{x}$-এর সংজ্ঞা থেকে
  $\Value{x}{M}[\Subst{s}{\Value{t'}{M}[s]}{x}] = \Value{t'}{M}[s]$।

\item $t \ident \Atom{f}{t_1,\dots,t_n}$ হলে পাই:
\begin{multline*}
  \Value{\Subst{t}{t'}{x}}{M}[s]  = \\
  \begin{aligned}[b]
& = \Value{\Atom{f}{\Subst{t_1}{t'}{x}, \dots, \Subst{t_n}{t'}{x}}}{M}[s]\\
& \qquad    \text{$\Subst{t}{t'}{x}$-এর সংজ্ঞা থেকে}\\
& = \Assign{f}{M}(\Value{\Subst{t_1}{t'}{x}}{M}[s], \dots,
    \Value{\Subst{t_n}{t'}{x}}{M}[s])\\
    & \qquad  \text{$\Value{\Atom{f}{\dots}}{M}[s]$-এর সংজ্ঞা থেকে}\\
& = \Assign{f}{M}(\Value{t_1}{M}[\Subst{s}{\Value{t'}{M}[s]}{x}], \dots,
   \Value{t_n}{M}[\Subst{s}{\Value{t'}{M}[s]}{x}])\\
& \qquad    \text{আরোহ-অনুমান থেকে}\\
& = \Value{t}{M}[\Subst{s}{\Value{t'}{M}[s]}{x}]
    \text{$\Value{\Atom{f}{\dots}}{M}[\Subst{s}{\Value{t'}{M}[s]}{x}]$-এর সংজ্ঞা থেকে}
  \end{aligned}
\end{multline*}
\end{enumerate}
\end{proof}

\begin{prop}\ollabel{prop:ext-formulas} ধরি $\Struct M$ একটি
!!a{structure}, $!A$ একটি !!a{formula}, $t'$ এমন একটি পদ যা~$!A$-তে
$x$-এর জন্য মুক্ত, এবং $s$ একটি চলরাশি-আরোপ। তাহলে
$\Sat{M}{\Subst{!A}{t'}{x}}[s]$ তখনই, যখন
$\Sat{M}{!A}[\Subst{s}{\Value{t'}{M}[s]}{x}]$।
\end{prop}
% BN-SRC-102 TARGET-CORRECTION-BEGIN
\par\smallskip\noindent\textnormal{\small\textbf{উৎস-সংশোধন (BN-SRC-102):} হিমায়িত ইংরেজি উৎসে সূত্র-প্রতিস্থাপন প্রতিজ্ঞায় $t'$-কে শুধু পদ বলা হয়েছে। কিন্তু এই অধ্যায়ের সূত্রে প্রতিস্থাপনের সংজ্ঞা $\Subst{!A}{t'}{x}$ কেবল তখনই দেয়, যখন $t'$,~$!A$-তে $x$-এর জন্য মুক্ত; চলরাশি-আবদ্ধতা এড়াতে এই অনুমান অপরিহার্য। বাংলা প্রতিজ্ঞায় তাই অনুমানটি স্পষ্টভাবে রাখা হয়েছে।} % BN-SRC-102 TARGET-CORRECTION-END

\begin{proof}
অনুশীলনী।
\end{proof}

\begin{prob}
\olref[fol][syn][ext]{prop:ext-formulas} প্রমাণ করুন।
\end{prob}

\begin{explain}
  \cref{fol:syn:ext:prop:ext-terms,fol:syn:ext:prop:ext-formulas}-এর
  বক্তব্য হলো এই। ধরি আমাদের একটি পদ $t$ অথবা !!a{formula}~$!A$ এবং
  কোনো পদ~$t'$ আছে; আমরা $\Subst{t}{t'}{x}$-এর মান, অথবা
  !!a{structure}~$\Struct M$-এ !!a{variable} আরোপ~$s$-এর সাপেক্ষে
  $\Subst{!A}{t'}{x}$ পরিতৃপ্ত কি না জানতে চাই। তাহলে হয় প্রথমে
  প্রতিস্থাপন করে পরে $\Struct{M}$ ও~$s$-এর সাপেক্ষে মান বা পরিতৃপ্তি
  বিবেচনা করতে পারি; অথবা প্রথমে $\Struct{M}$-এ $s$-এর সাপেক্ষে
  $t'$-এর মান~$m = \Value{t'}{M}[s]$ নির্ণয় করে !!{variable} আরোপকে
  $\Subst{s}{m}{x}$-এ বদলে তারপর $\Struct{M}$ ও~$\Subst{s}{m}{x}$-এর
  সাপেক্ষে~$t$-এর মান, অথবা $\Sat{M}{!A}[\Subst{s}{m}{x}]$ কি না
  বিবেচনা করতে পারি।
  \Cref{fol:syn:ext:prop:ext-terms,fol:syn:ext:prop:ext-formulas}
  নিশ্চিত করে, যে পথেই করি না কেন উত্তর একই হবে।
\end{explain}

\end{document}
